
% appendices/appendix_d_references.tex

\chapter{参考资料与延伸阅读}

AI Infrastructure 是一个跨越算法、系统、硬件、网络和平台工程的领域。单靠一本书无法覆盖所有细节。本附录按照主题列出建议继续阅读的论文、系统、文档和开源项目。阅读这些资料时，建议不要只关注结论，而要关注它们解决的问题、隐含的硬件假设、系统瓶颈和工程取舍。

\section{Transformer 与 LLM 论文}

\subsection{基础架构}

\begin{itemize}
\item Vaswani et al., \textit{Attention Is All You Need}. Transformer 架构的奠基论文，建议重点阅读 Scaled Dot-Product Attention、Multi-Head Attention 和 Position-wise FFN 的设计动机。
\item Radford et al., \textit{Improving Language Understanding by Generative Pre-Training}. GPT 系列早期工作，展示了预训练语言模型的基本范式。
\item Brown et al., \textit{Language Models are Few-Shot Learners}. GPT-3 论文，建议关注规模扩展、in-context learning 和训练数据混合方式。
\item Touvron et al., \textit{LLaMA: Open and Efficient Foundation Language Models}. 理解现代开源 Decoder-only LLM 的重要参考。
\item Touvron et al., \textit{Llama 2: Open Foundation and Fine-Tuned Chat Models}. 建议结合第 6 章阅读，关注预训练、SFT 和对齐流程。
\end{itemize}

\subsection{位置编码与结构改进}

\begin{itemize}
\item Su et al., \textit{RoFormer: Enhanced Transformer with Rotary Position Embedding}. RoPE 的代表性论文。
\item Shazeer, \textit{GLU Variants Improve Transformer}. 理解 SwiGLU、GEGLU 等 FFN 变体的参考。
\item Shazeer, \textit{Fast Transformer Decoding: One Write-Head is All You Need}. Multi-Query Attention 的重要参考。
\item Ainslie et al., \textit{GQA: Training Generalized Multi-Query Transformer Models from Multi-Head Checkpoints}. Grouped-Query Attention 的重要参考。
\end{itemize}

\subsection{规模定律与训练范式}

\begin{itemize}
\item Kaplan et al., \textit{Scaling Laws for Neural Language Models}. 早期语言模型规模定律研究。
\item Hoffmann et al., \textit{Training Compute-Optimal Large Language Models}. Chinchilla Scaling Law 的代表性工作。
\item Ouyang et al., \textit{Training language models to follow instructions with human feedback}. InstructGPT 与 RLHF 的重要参考。
\end{itemize}

\section{分布式训练论文}

\subsection{数据并行与集合通信}

\begin{itemize}
\item Li et al., \textit{PyTorch Distributed: Experiences on Accelerating Data Parallel Training}. 理解 PyTorch DDP 工程设计的重要资料。
\item Sergeev and Del Balso, \textit{Horovod: Fast and Easy Distributed Deep Learning in TensorFlow}. Ring AllReduce 训练实践的重要系统。
\item NVIDIA NCCL Documentation. 理解 AllReduce、AllGather、ReduceScatter 等通信原语及其性能调优。
\end{itemize}

\subsection{模型并行}

\begin{itemize}
\item Shoeybi et al., \textit{Megatron-LM: Training Multi-Billion Parameter Language Models Using Model Parallelism}. Tensor Parallelism 和大模型训练的重要论文。
\item Narayanan et al., \textit{Efficient Large-Scale Language Model Training on GPU Clusters Using Megatron-LM}. 3D 并行、流水线调度和大规模训练实践的重要参考。
\item Huang et al., \textit{GPipe: Efficient Training of Giant Neural Networks using Pipeline Parallelism}. Pipeline Parallelism 的经典论文。
\item Narayanan et al., \textit{PipeDream: Generalized Pipeline Parallelism for DNN Training}. 理解流水线并行调度的重要参考。
\end{itemize}

\subsection{ZeRO 与内存优化}

\begin{itemize}
\item Rajbhandari et al., \textit{ZeRO: Memory Optimizations Toward Training Trillion Parameter Models}. ZeRO Stage 1/2/3 的核心论文。
\item Rajbhandari et al., \textit{ZeRO-Offload: Democratizing Billion-Scale Model Training}. CPU offload 的重要参考。
\item Ren et al., \textit{ZeRO-Infinity: Breaking the GPU Memory Wall for Extreme Scale Deep Learning}. CPU / NVMe offload 与更大模型训练的参考。
\item Chen et al., \textit{Training Deep Nets with Sublinear Memory Cost}. Gradient Checkpointing 的经典工作。
\end{itemize}

\section{Serving Infra 论文与系统}

\subsection{KV Cache 与推理调度}

\begin{itemize}
\item Kwon et al., \textit{Efficient Memory Management for Large Language Model Serving with PagedAttention}. vLLM 与 PagedAttention 的核心论文。
\item Yu et al., \textit{Orca: A Distributed Serving System for Transformer-Based Generative Models}. Continuous Batching 和 iteration-level scheduling 的重要系统。
\item NVIDIA FasterTransformer / TensorRT-LLM Documentation. 理解高性能 LLM 推理 kernel、engine 构建和部署优化的重要资料。
\end{itemize}

\subsection{Attention 优化}

\begin{itemize}
\item Dao et al., \textit{FlashAttention: Fast and Memory-Efficient Exact Attention with IO-Awareness}. FlashAttention 的核心论文。
\item Dao, \textit{FlashAttention-2: Faster Attention with Better Parallelism and Work Partitioning}. 理解更高效 attention kernel 并行划分的参考。
\item Shah et al., \textit{FlashAttention-3: Fast and Accurate Attention with Asynchrony and Low-precision}. 理解新一代 GPU 上 attention 优化方向的参考。
\end{itemize}

\subsection{量化与压缩}

\begin{itemize}
\item Frantar et al., \textit{GPTQ: Accurate Post-Training Quantization for Generative Pre-trained Transformers}. GPTQ 的代表性论文。
\item Lin et al., \textit{AWQ: Activation-aware Weight Quantization for LLM Compression and Acceleration}. AWQ 的代表性论文。
\item Dettmers et al., \textit{LLM.int8(): 8-bit Matrix Multiplication for Transformers at Scale}. 大模型 INT8 推理的重要参考。
\item Dettmers et al., \textit{QLoRA: Efficient Finetuning of Quantized LLMs}. 量化与参数高效微调结合的重要参考。
\end{itemize}

\section{CUDA 与 GPU 优化资料}

\subsection{官方文档}

\begin{itemize}
\item NVIDIA CUDA C++ Programming Guide. 学习 CUDA 编程模型、线程层次、内存模型和同步机制的基础资料。
\item NVIDIA CUDA C++ Best Practices Guide. 学习 coalescing、shared memory、occupancy、异步拷贝等优化技巧。
\item NVIDIA Nsight Systems Documentation. 分析端到端 timeline、CPU-GPU 协作和通信重叠。
\item NVIDIA Nsight Compute Documentation. 分析单个 CUDA kernel 的 occupancy、memory throughput、warp stall 和 Tensor Core 利用率。
\end{itemize}

\subsection{矩阵计算与 Kernel 优化}

\begin{itemize}
\item CUTLASS Documentation. 理解 tile-based GEMM、Tensor Core MMA 和高性能矩阵乘法模板的重要资料。
\item Triton Documentation. 学习使用 Python-like DSL 编写高性能 GPU kernel。
\item cuBLAS Documentation. 理解 GEMM、batched GEMM 和 mixed precision 矩阵计算。
\item cuDNN Documentation. 理解深度学习基础算子的高性能实现。
\end{itemize}

\section{网络与集群调度资料}

\subsection{高性能网络}

\begin{itemize}
\item InfiniBand Architecture Specification. 理解 InfiniBand 网络、queue pair、completion queue 和 RDMA 语义的基础资料。
\item NVIDIA GPUDirect RDMA Documentation. 理解 GPU memory 与 NIC 直接通信的关键资料。
\item RoCE Deployment Guides. 理解 RoCEv2、PFC、ECN、lossless Ethernet 与拥塞控制。
\item NCCL Tuning Guide. 理解多 NIC、多 rail、拓扑感知通信和环境变量调优。
\end{itemize}

\subsection{Kubernetes 与 AI 调度}

\begin{itemize}
\item Kubernetes Documentation. Pod、Node、Scheduler、Device Plugin、Operator 等基础概念。
\item NVIDIA Kubernetes Device Plugin. GPU 资源暴露、调度与容器运行时集成的重要项目。
\item Volcano Documentation. Gang Scheduling 和 AI/HPC batch workload 调度的重要系统。
\item Kueue Documentation. Kubernetes 原生队列管理、配额和批作业调度的重要项目。
\item Kubeflow Training Operator Documentation. 在 Kubernetes 上运行分布式训练任务的重要参考。
\end{itemize}

\section{开源项目索引}

\subsection{训练框架与并行训练}

\begin{longtable}{p{0.26\textwidth}p{0.64\textwidth}}
\toprule
项目 & 关注点 \
\midrule
PyTorch Distributed & DDP、ProcessGroup、分布式通信基础。 \
DeepSpeed & ZeRO、offload、混合精度、大规模训练优化。 \
Megatron-LM & Tensor Parallelism、Pipeline Parallelism、3D Parallelism。 \
Hugging Face Accelerate & 多设备训练启动与配置抽象。 \
Colossal-AI & 并行训练、内存优化和大模型训练工具链。 \
FairScale & FSDP、checkpointing 等分布式训练能力。 \
\bottomrule
\end{longtable}

\subsection{推理服务}

\begin{longtable}{p{0.26\textwidth}p{0.64\textwidth}}
\toprule
项目 & 关注点 \
\midrule
vLLM & PagedAttention、Continuous Batching、高吞吐 LLM Serving。 \
TensorRT-LLM & NVIDIA GPU 上的高性能 LLM 推理优化。 \
Text Generation Inference & 面向生产部署的文本生成推理服务。 \
SGLang & LLM serving runtime、structured generation 与调度优化。 \
llama.cpp & CPU / GPU 本地推理、量化格式和轻量部署。 \
ONNX Runtime & 通用模型推理优化和跨硬件执行。 \
\bottomrule
\end{longtable}

\subsection{可观测性与运维}

\begin{longtable}{p{0.26\textwidth}p{0.64\textwidth}}
\toprule
项目 & 关注点 \
\midrule
Prometheus & 指标采集与时序存储。 \
Grafana & 监控面板与可视化。 \
DCGM Exporter & NVIDIA GPU 指标采集。 \
OpenTelemetry & 分布式 tracing、metrics、logs 统一观测。 \
Loki & 日志聚合。 \
Jaeger & 分布式链路追踪。 \
\bottomrule
\end{longtable}

\section{阅读建议}

面对大量资料，建议按问题驱动阅读，而不是按论文发布时间线机械阅读：

\begin{enumerate}
\item 如果目标是理解模型结构，先读 Transformer、GPT、LLaMA、RoPE、GQA 等资料。
\item 如果目标是训练大模型，先读 DDP、Megatron-LM、ZeRO、GPipe、NCCL 文档。
\item 如果目标是优化推理服务，先读 KV Cache、PagedAttention、FlashAttention、Continuous Batching 和量化论文。
\item 如果目标是做 GPU kernel 优化，先读 CUDA Programming Guide、Best Practices Guide、CUTLASS 和 Triton 文档。
\item 如果目标是建设 AI 集群，先读 NCCL、GPUDirect RDMA、Kubernetes Device Plugin、Volcano、Kueue 和监控系统资料。
\end{enumerate}

读论文时可以重点问五个问题：

\begin{itemize}
\item 它解决的瓶颈是什么：计算、显存、通信、调度还是工程复杂度？
\item 它假设的硬件环境是什么：单卡、多卡、NVLink、RDMA、HBM 容量？
\item 它优化的是训练还是推理，是吞吐还是延迟？
\item 它引入了哪些新的 trade-off？
\item 它在真实生产系统中落地时还缺什么？
\end{itemize}

\section{附录 D 小结}

AI Infrastructure 的知识更新很快，但底层问题相对稳定：如何让计算更满、显存更省、通信更少、调度更准、故障更可控。持续阅读论文和开源系统时，应始终把注意力放在这些长期问题上，而不是只追逐某个短期流行框架或参数配置。
